\section{Additional tables}\label{app:tables}

\begin{table}[h]
\centering
\tabRealization
\caption{Net revenue under Rapid~/~proportional by realized share of the
incremental capital flow, \$B, 2030.}
\label{tab:realization}
\end{table}

\begin{table}[h]
\centering
\tabCapitalScope
\caption{Capital-scope sensitivity: net revenue with and without taxable
pension and IRA distributions in the capital set (proportional variants),
\$B, 2030.}
\label{tab:scope}
\end{table}

\begin{table}[h]
\centering
\tabStates
\caption{Top ten states by net state revenue change (state tax before
refundable credits minus refundable state credits), Rapid~/~proportional,
\$B, 2030.}
\label{tab:states}
\end{table}

\begin{table}[h]
\centering
\tabSweep
\caption{The mechanism sweep, \genSweepYear{}: net revenue, federal income
tax, SPM poverty and net Gini by share of labor income shifted to capital.}
\label{tab:sweep}
\end{table}

\clearpage
\section{Shock routing into the tax base}\label{app:routing}

Several shocked variables are aggregates rather than raw inputs, so the
implementation verifies that overriding them moves taxable income rather
than being recomputed away. Traced against the model's gross-income and
market-income source registries (\texttt{gov.irs.gross\_income.sources} and
\texttt{gov.household.market\_income\_sources} in policyengine-us
\genModelVersion{}):

\begin{table}[h]
\centering
\small
\begin{tabular}{lll}
\toprule
Shocked variable & Route into IRS gross income & Route into market income \\
\midrule
employment income & via employment income, net of & listed source \\
 & pre-tax contributions & \\
self-employment income & listed source & listed source \\
long/short-term capital gains & capital gains & capital gains \\
qualified/non-qualified dividends & ordinary dividend income & dividend income \\
taxable interest & listed source & interest income \\
rental income & listed source & listed source \\
partnership and S-corp income & listed source & listed source \\
taxable pension income & listed source & pension income \\
taxable IRA distributions & taxable retirement distributions & retirement distributions \\
\bottomrule
\end{tabular}
\caption{Every shocked variable reaches both the federal tax base and
household market income. Qualified and non-qualified dividends are the two
components of ordinary dividends; short- and long-term gains the two
components of capital gains; tax-exempt interest is deliberately outside the
shocked set.}
\label{tab:routing}
\end{table}

Because every shocked variable reaches household market income, the
net-income identity closes; the maximum residual across scenarios is
\$\genIdentityMaxResidual B, and the labor aggregate lands on its target to
within relative error \genLaborTargetMaxErr{}.

\clearpage
\section{Reproducibility}\label{app:repro}

Everything in this paper runs from open source against openly published
microdata:

\begin{itemize}
\item \textbf{Model:} policyengine-us \genModelVersion{} via policyengine.py
\genRunnerVersion{} (the runtime certifies the model--data pairing).
\item \textbf{Data:} \texttt{populace-us} certified build
\texttt{\genDataBuild}.
\item \textbf{Code:} \texttt{github.com/PolicyEngine/ai-inequality} ---
scenario engine in \texttt{analysis/ai\_scenarios.py}, runner in
\texttt{analysis/compute\_ai\_scenarios.py}, sweep in
\texttt{analysis/compute\_shift\_sweep.py}, reconciliation in
\texttt{analysis/reconcile\_budget\_lab.py}. Every number and table in this
paper is emitted from the run outputs by
\texttt{analysis/emit\_paper\_values.py}; figures by
\texttt{analysis/paper\_figures.py}.
\item \textbf{Their side:} the Budget Lab's \texttt{AI-Fiscal} repository
(MIT license) commits the complete published result grid
(\texttt{results/aggregates/ai\_fiscal\_publishable\_2030\_latest.xlsx});
the reconciliation reads it directly, and a copy is committed alongside this
paper's own run outputs so the bridge is reproducible without a network
fetch. Their PUF-based input microdata cannot be redistributed, so their
pipeline is readable but not rerunnable by third parties; this paper's is
both.
\item \textbf{Canonical run outputs:} \texttt{analysis/outputs/ai\_scenarios.json}
(build~p) is committed with three comparison files:
\texttt{ai\_scenarios\_buildp\_published.json} (build~p as first published,
with Head Start and Early Head Start in net income) and
\texttt{ai\_scenarios\_buildo.json} (build~o), the stability comparison of
Section~\ref{sec:stability}; and \texttt{transfer\_detail\_buildp\_published.json},
program-by-program benefit totals on the published runtime, the source of the
Head Start amounts in the correction note. Every value and table in this
paper regenerates from the repository alone --- the full simulation need only
be rerun to change a scenario, not to check a number.
\item \textbf{Interactive companion:}
\url{https://policyengine.org/us/ai-inequality/income-shift} (scenario
explorer) and \url{https://policyengine.org/us/ai-inequality/budget-lab}
(memo-length summary).
\end{itemize}

Calibration unit tests reproduce all nine of the Budget Lab's published
derived scenario values to within 0.005\,pp from the two recovered
constants ($G_{\mathrm{CBO}}=1.0976$, $\theta^0_L=0.5550$) and re-derive the
labor share by grid search. The scenario runner checkpoints per-scenario
results and is resumable; run logs and build-o archives are kept under
\texttt{analysis/outputs/}.
